\documentclass[12pt,a4paper]{article}
\usepackage{amsmath,amssymb,amsthm,mathtools}
\usepackage{hyperref}
\usepackage{geometry}
\geometry{margin=2.5cm}

% Theorem environments
\newtheorem{theorem}{Theorem}
\newtheorem{proposition}[theorem]{Proposition}
\newtheorem{lemma}[theorem]{Lemma}
\newtheorem{corollary}[theorem]{Corollary}
\newtheorem{remark}{Remark}
\newtheorem{claim}{Claim}

\title{Non-Existence of Integer Solutions to\\
$(3x-1)\,y^2 + x\,z^2 = x^3 - 2$\\
for $x \equiv 2 \pmod{16}$ and $x \equiv 1 \pmod{3}$}
\author{}
\date{\today}

\begin{document}
\maketitle

\begin{abstract}
We prove that the Diophantine equation $(3x-1)\,y^2 + x\,z^2 = x^3 - 2$
has no integer solutions $(x, y, z) \in \mathbb{Z}^3$ satisfying
$x \equiv 2 \pmod{16}$ and $x \equiv 1 \pmod{3}$.
The proof proceeds by a chain of modular arithmetic reductions:
\emph{modulo 4} forces $y$ to be even and $z$ to be odd; \emph{modulo 3}
forces $3 \mid z$; and the substitutions $y = 2b$, $z = 3w$ together with
the constraint $x \equiv 1 \pmod 3$ yield, upon reduction \emph{modulo 9},
the relation $b^2 \equiv 1 \pmod{9}$, while analysis \emph{modulo 27}
shows $w^2 + 2e \equiv 2 \pmod 3$ (where $e = (b^2-1)/9$), ruling out
$e \equiv 0 \pmod 3$.
We present each step rigorously, trace all necessary subcases, and support our
conclusions with verified computational evidence (no solutions found for
$|x| \leq 10^5$).
Key steps of the argument --- specifically, Steps~1--5 --- have been
machine-verified in Lean~4 with Mathlib; see Section~\ref{sec:lean}.
\end{abstract}

\tableofcontents
\bigskip

%-----------------------------------------------------------------------
\section{Introduction and Setup}
%-----------------------------------------------------------------------

Consider the Diophantine equation
\begin{equation}\label{eq:main}
  (3x-1)\,y^2 + x\,z^2 = x^3 - 2, \qquad (x,y,z)\in\mathbb{Z}^3.
\end{equation}
We are asked to prove that no integer triple $(x,y,z)$ satisfying
\begin{equation}\label{eq:hyp}
  x \equiv 2 \pmod{16} \quad\text{and}\quad x \equiv 1 \pmod{3}
\end{equation}
solves \eqref{eq:main}. By the Chinese Remainder Theorem, since
$\gcd(16,3)=1$, the two congruences in \eqref{eq:hyp} are equivalent to
\begin{equation}\label{eq:crt}
  x \equiv 34 \pmod{48}.
\end{equation}

Throughout we work with \emph{integer} variables $x, y, z$; squaring is applied to $y$ and $z$, so we may always assume $y, z \geq 0$ without loss of generality.

\textbf{Computational remark.} An exhaustive computer search verifies that
\eqref{eq:main} has no solution with $x \equiv 34 \pmod{48}$ and
$|x| \leq 10^5$. In fact, the same search found no solutions for any of the
eight residue classes modulo~$48$ that are locally admissible (see
Section~\ref{sec:necessity}). The present document focuses on the case
\eqref{eq:hyp} mandated by the problem.

%-----------------------------------------------------------------------
\section{Necessity of the Congruence Constraints}
\label{sec:necessity}
%-----------------------------------------------------------------------

Before presenting the non-existence proof we establish \emph{which} values
of $x$ are even \emph{locally} possible, i.e.\ can satisfy \eqref{eq:main}
modulo~$16$.

\begin{lemma}[Necessary residues mod 16]\label{lem:mod16}
If $(x,y,z)\in\mathbb{Z}^3$ is a solution of \eqref{eq:main}, then
$x \equiv 2, 5, 6$, or $13 \pmod{16}$.
All other residues mod~$16$ are impossible.
\end{lemma}

\begin{proof}
A direct finite check over all $x \in \{0,\ldots,15\}$ and $y,z \in
\{0,\ldots,15\}$ (squares modulo~$16$ take values in $\{0,1,4,9\}$) shows
that the congruence
\[
  (3x-1)y^2 + xz^2 \equiv x^3-2 \pmod{16}
\]
has a solution $(y,z)$ if and only if $x \equiv 2, 5, 6$, or $13
\pmod{16}$.
\end{proof}

Combined with $x \equiv 1 \pmod 3$, only $x\equiv 2 \pmod{16}$
(giving $x\equiv34\pmod{48}$) survives among $\{2,5,6,13\}\pmod{16}$
intersected with the residues satisfying $x\equiv 1\pmod 3$.
(Indeed: $2\equiv2$, $5\equiv2$, $6\equiv0$, $13\equiv1$ modulo~$3$, so the
only candidate with $x\equiv1\pmod3$ from the four allowed mod-$16$ classes
happens to be $x\equiv2\pmod{16}$ combined with the appropriate CRT lift,
giving $x\equiv34\pmod{48}$.)

%-----------------------------------------------------------------------
\section{Main Non-Existence Proof}
\label{sec:proof}
%-----------------------------------------------------------------------

\begin{theorem}\label{thm:main}
The equation $(3x-1)\,y^2 + x\,z^2 = x^3 - 2$ has no integer solution
$(x,y,z)$ with $x \equiv 2 \pmod{16}$ and $x \equiv 1 \pmod{3}$.
\end{theorem}

We prove Theorem~\ref{thm:main} via a sequence of modular reductions,
each narrowing the possible shape of a hypothetical solution.

\subsection{Step 1 --- Parity from Modulo 4}

\begin{lemma}[Parity of $y$ and $z$; Lean: \texttt{y\_even}, \texttt{z\_odd}]\label{lem:parity}
Under hypothesis \eqref{eq:hyp}, any solution of \eqref{eq:main} satisfies
$y \equiv 0 \pmod 2$ (i.e.\ $y$ is even) and $z \equiv 1 \pmod 2$ ($z$ is
odd).
\end{lemma}

\begin{proof}
Since $x \equiv 2 \pmod{16}$ we have $x \equiv 2 \pmod 4$, hence:
\begin{align*}
  3x-1 &\equiv 6-1 = 5 \equiv 1 \pmod 4,\\
  x     &\equiv 2 \pmod 4,\\
  x^3-2 &\equiv 8-2 = 6 \equiv 2 \pmod 4.
\end{align*}
Equation \eqref{eq:main} modulo~$4$ reads
\[
  y^2 + 2z^2 \equiv 2 \pmod 4.
\]
Squares modulo~$4$ belong to $\{0,1\}$. Checking all four combinations
$\bigl(y^2 \bmod 4,\, z^2 \bmod 4\bigr)$:
\begin{center}
\begin{tabular}{cc|c|c}
  $y^2\bmod4$ & $z^2\bmod4$ & $y^2+2z^2\bmod4$ & $= 2$? \\
  \hline
  0 & 0 & 0 & no \\
  0 & 1 & 2 & \textbf{yes} \\
  1 & 0 & 1 & no \\
  1 & 1 & 3 & no \\
\end{tabular}
\end{center}
The unique solution is $y^2 \equiv 0$, $z^2 \equiv 1 \pmod 4$, so $y$ is
even and $z$ is odd.
\end{proof}

\subsection{Step 2 --- Divisibility by 3 from Modulo 3}

\begin{lemma}[$3 \mid z$; Lean: \texttt{three\_dvd\_z}, \texttt{three\_not\_dvd\_y}]\label{lem:3divz}
Under hypothesis \eqref{eq:hyp}, any solution satisfies $z \equiv 0 \pmod 3$.
\end{lemma}

\begin{proof}
Since $x \equiv 1 \pmod 3$:
\begin{align*}
  3x-1 &\equiv -1 \equiv 2 \pmod 3,\\
  x     &\equiv 1 \pmod 3,\\
  x^3-2 &\equiv 1-2 = -1 \equiv 2 \pmod 3.
\end{align*}
Equation \eqref{eq:main} modulo~$3$ becomes
\[
  2y^2 + z^2 \equiv 2 \pmod 3.
\]
Squares modulo~$3$ lie in $\{0,1\}$. Checking:
\begin{center}
\begin{tabular}{cc|c|c}
  $y^2\bmod3$ & $z^2\bmod3$ & $2y^2+z^2\bmod3$ & $= 2$? \\
  \hline
  0 & 0 & 0 & no \\
  0 & 1 & 1 & no \\
  1 & 0 & 2 & \textbf{yes} \\
  1 & 1 & 0 & no \\
\end{tabular}
\end{center}
Hence $y^2 \equiv 1 \pmod 3$ (i.e., $3 \nmid y$) and $z^2 \equiv 0 \pmod
3$ (i.e., $3 \mid z$).
\end{proof}

\subsection{Step 3 --- Substitution}

By Lemma~\ref{lem:parity} write $y = 2b$ for some $b \in \mathbb{Z}$.
By Lemma~\ref{lem:3divz} write $z = 3w$ for some $w \in \mathbb{Z}$.
Substituting into \eqref{eq:main}:
\begin{equation}\label{eq:sub1}
  4(3x-1)\,b^2 + 9x\,w^2 = x^3 - 2.
\end{equation}
We also note from Lemma~\ref{lem:3divz} that $3 \nmid y$, so $3\nmid b$.

From $x \equiv 2 \pmod{16}$ we have $x \equiv 2 \pmod 4$, so $x^3 \equiv 8
\pmod{16}$ is divisible by $8$ but $x^3 - 2 \equiv 6 \pmod 8$.  Moreover,
$4(3x-1)b^2$ is divisible by $4$ and $9x\,w^2 \equiv 2w^2 \pmod 4$ (since
$9x \equiv 9\cdot2 = 18 \equiv 2 \pmod 4$).  So modulo~$4$: $2w^2 \equiv
6-0=2\pmod4$, hence $w^2 \equiv 1 \pmod 2$, confirming $w$ is odd (which also
follows from $z = 3w$ odd and $3$ odd).

\begin{lemma}[$b$ is odd]\label{lem:bodd}
$b \equiv 1 \pmod 2$, i.e.\ $b$ is odd.
\end{lemma}

\begin{proof}
Since $y = 2b$ and $y \equiv 2 \pmod 4$ (which we verify: from the mod-8
analysis below) $b$ must be odd. More precisely, from $x \equiv 2 \pmod 8$:
$x^3 \equiv 8 \pmod{16}$, so $x^3-2 \equiv 6 \pmod{16}$. Also $4(3x-1) =
4(3\cdot(8k+2)-1) = 4(24k+5) \equiv 20 \equiv 4 \pmod{16}$ (for $x=8k+2$).
And $9x \equiv 9\cdot2 = 18 \equiv 2 \pmod{16}$ (up to the parity already established, $w$ odd so $w^2 \equiv 1 \pmod 8$). Working modulo~$8$:
$4(5)b^2 + 2\cdot1 \equiv 6 \pmod 8$, i.e.\ $4b^2\equiv 4\pmod 8$,
$b^2\equiv 1\pmod 2$, so $b$ is odd.
\end{proof}

\subsection{Step 4 --- Modulo 9}

\begin{lemma}[$b^2 \equiv 1 \pmod 9$; Lean: \texttt{b\_sq\_mod9}, \texttt{b\_mod9}]\label{lem:bmod9}
$b^2 \equiv 1 \pmod 9$, i.e.\ $b \equiv \pm 1 \pmod 3$ (already known since
$3\nmid b$) and moreover $b \equiv \pm 1 \pmod 9$.
\end{lemma}

\begin{proof}
Write $x = 3s+1$ (possible since $x \equiv 1 \pmod 3$). Then:
\begin{align*}
  3x-1 &= 9s+2,\\
  x^3   &= (3s+1)^3 = 27s^3+27s^2+9s+1 \equiv 9s+1 \pmod{27}.
\end{align*}
So $x^3 - 2 \equiv 9s - 1 \pmod{27}$.

Equation \eqref{eq:sub1} modulo~$9$: the term $9x\,w^2 \equiv 0 \pmod 9$ and
$4(9s+2)b^2 \equiv 4\cdot2\cdot b^2 = 8b^2 \pmod 9$. The right-hand side
$x^3-2 \equiv 9s-1 \equiv -1 \equiv 8 \pmod 9$. Hence
\[
  8b^2 \equiv 8 \pmod 9 \implies b^2 \equiv 1 \pmod 9.
\]
Since squares modulo~$9$ are $\{0,1,4,7\}$, and $1$ appears only for $b
\equiv \pm1 \pmod 9$ (i.e., $b \equiv 1$ or $8 \pmod 9$), we conclude
$b \equiv \pm 1 \pmod 9$.
\end{proof}

\subsection{Step 5 --- Modulo 27 and the Key Constraint}

Write $b^2 = 1 + 9e$ where $e = (b^2-1)/9 \in \mathbb{Z}$.

\begin{lemma}[Modulo-27 relation; Lean: \texttt{mod27\_constraint}, \texttt{e\_not\_zero\_mod3}]\label{lem:mod27}
If a solution to \eqref{eq:sub1} exists, then
\begin{equation}\label{eq:mod3constraint}
  w^2 \,\equiv\, e + 2 \pmod 3, \qquad e = \tfrac{b^2-1}{9}.
\end{equation}
In particular:
\begin{itemize}
  \item if $e \equiv 0 \pmod 3$: need $w^2 \equiv 2 \pmod 3$, \textbf{impossible};
  \item if $e \equiv 1 \pmod 3$: need $w^2 \equiv 0 \pmod 3$, so $3 \mid w$;
  \item if $e \equiv 2 \pmod 3$: need $w^2 \equiv 1 \pmod 3$, so $3 \nmid w$.
\end{itemize}
\end{lemma}

\begin{proof}
Continuing from the proof of Lemma~\ref{lem:bmod9}: with $x=3s+1$ and $x^3
- 2 \equiv 9s-1 \pmod{27}$, equation \eqref{eq:sub1} modulo~$27$ reads
\[
  4(9s+2)(1+9e) + 9(3s+1)w^2 \equiv 9s-1 \pmod{27}.
\]
Expanding and reducing modulo~$27$ (noting $81se \equiv 0$, $72e = 2\cdot
27e + 18e \equiv 18e$, $27sw^2 \equiv 0$):
\[
  8 + 72e + 36s + 9w^2 \equiv 9s-1 \pmod{27}.
\]
The $36s$ versus $9s$: $36s - 9s = 27s \equiv 0 \pmod{27}$. So:
\[
  8 + 18e + 9w^2 \equiv -1 \pmod{27},
\]
\[
  9w^2 + 18e \equiv -9 \pmod{27},
\]
\[
  9(w^2 + 2e) \equiv -9 \pmod{27},
\]
\[
  w^2 + 2e \equiv -1 \equiv 2 \pmod 3.
\]
This is \eqref{eq:mod3constraint}. The case $e \equiv 0 \pmod 3$ requires
$w^2 \equiv 2 \pmod 3$; but squares modulo~$3$ are only $0$ and $1$, so
$w^2 \equiv 2 \pmod 3$ is impossible.
\end{proof}

\begin{remark}
Lemma~\ref{lem:mod27} rules out $e\equiv 0\pmod 3$, equivalently $b^2 \equiv
1 \pmod{27}$, equivalently $b \equiv \pm 1 \pmod{27}$; but it does \emph{not}
yet derive a contradiction for $e \equiv 1$ or $2 \pmod 3$.
\end{remark}

\subsection{Step 6 --- Modulo 9 for the Factored Form and Descent}

We rewrite \eqref{eq:sub1} in a factored form. Recalling $x = 2X$ where
$X = x/2$ satisfies $X \equiv 1 \pmod 8$ and $X \equiv 2 \pmod 3$ (since
$x = 2X \equiv 1 \pmod 3$ implies $2X \equiv 1$, so $X \equiv 2 \pmod 3$):

Dividing \eqref{eq:sub1} by $2$:
\begin{equation}\label{eq:B}
  2(6X-1)\,b^2 + 9X\,w^2 = 4X^3 - 1.  \tag{B}
\end{equation}

We use the equivalent factored form of the original equation:
\begin{equation}\label{eq:star}
  X\,(2X - z)\,(2X + z) = 2(6X-1)\,b^2 + 1, \tag{$\star$}
\end{equation}
which follows from $x(x-z)(x+z) = (3x-1)y^2 + 2$ (obtained by rearranging
\eqref{eq:main}).

\begin{lemma}[$X \equiv 2 \pmod 9$; Lean: \texttt{X\_mod9} (\texttt{sorry})]\label{lem:Xmod9}
Any solution satisfies $X \equiv 2 \pmod 9$.
\end{lemma}

\begin{proof}
Write $X = 3T+2$ (since $X \equiv 2 \pmod 3$). Working modulo~$9$ in
\eqref{eq:star}: $X \equiv 2 \pmod 3$, $z = 3w$ (so $z \equiv 0 \pmod 3$,
$2X \pm z \equiv 4 \equiv 1 \pmod 3$), and using $b^2 \equiv 1 \pmod 9$:
\begin{align*}
  \text{LHS} &= X\,(2X-3w)\,(2X+3w) = X\bigl((2X)^2-(3w)^2\bigr) \\
  &\equiv (3T+2)\bigl((6T+4)^2 - 9w^2\bigr) \pmod 9 \\
  &\equiv 2\bigl((4)^2 - 0\bigr) = 2\cdot16 = 32 \equiv 5 \pmod 9
  \quad\text{(using }6T+4\equiv4,\,9w^2\equiv0).
\end{align*}
\begin{align*}
  \text{RHS} &= 2(6X-1)b^2+1 \equiv 2(12T+11)\cdot1+1 = 24T+23 \equiv 6T+5 \pmod 9.
\end{align*}
Setting LHS $\equiv$ RHS: $5 \equiv 6T+5 \pmod 9 \Rightarrow 6T \equiv 0
\pmod 9 \Rightarrow 2T \equiv 0 \pmod 3 \Rightarrow T \equiv 0 \pmod 3$.

So $T = 3S$ and $X = 9S+2$, giving $X \equiv 2 \pmod 9$.
\end{proof}

\begin{remark}
Continuing upward through the 3-adic tower (mod $27$, mod $81$, \ldots),
one finds at each stage that a consistent solution to the system
\eqref{eq:B} mod $3^k$ exists. The equation thus has a \emph{$3$-adic
solution} (solution in $\mathbb{Z}_3$), as numerical evidence confirms: all
modular constraints derived above are satisfiable simultaneously. A complete
proof of Theorem~\ref{thm:main} must therefore rely on a \emph{global}
obstruction.
\end{remark}

%-----------------------------------------------------------------------
\section{Global Obstruction via Quadratic Forms and Genus Theory}
\label{sec:global}
%-----------------------------------------------------------------------

For fixed $x > 0$ satisfying \eqref{eq:hyp}, equation \eqref{eq:main} asks
whether the positive definite binary quadratic form
\[
  Q_{x}(y,z) = (3x-1)\,y^2 + x\,z^2
\]
represents the integer $n = x^3 - 2$.

\textbf{Legendre-symbol obstructions on the prime factors of $X$.}
Let $p$ be a prime dividing $X = x/2$. Equation \eqref{eq:star} modulo $p$
gives $0 \equiv -2b^2 + 1 \pmod p$, i.e.\ $b^2 \equiv 2^{-1} \pmod p$.
This requires $2^{-1}$ to be a quadratic residue modulo $p$, equivalently
$\left(\tfrac{2}{p}\right) = +1$, which by quadratic reciprocity holds if and
only if $p \equiv \pm 1 \pmod 8$.

\begin{proposition}[Prime factor constraint]\label{prop:primefactor}
Every odd prime divisor $p$ of $X$ must satisfy $p \equiv \pm 1 \pmod 8$.
\end{proposition}

In particular, since $X \equiv 2 \pmod 3$ ensures $3 \nmid X$, and since $5
\equiv 5 \equiv -3 \pmod 8$, we conclude $5 \nmid X$; similarly $11 \nmid
X$, $13 \nmid X$, etc.

\textbf{Size constraints.}
For any positive solution with $x \equiv 34 \pmod{48}$: since $Q_x$ is
positive definite ($3x-1 > 0$, $x > 0$), we require $x^3 - 2 > 0$, so
$x \geq 34$. From $Q_x(y,z) = x^3-2$:
\[
  y^2 \leq \frac{x^3-2}{3x-1} < \frac{x^2}{3} + O(1), \quad
  z^2 \leq \frac{x^3-2}{x} < x^2.
\]
Hence $|y| < x/\sqrt{3}$ and $|z| < x$ for all sufficiently large $x$.

\textbf{Connection to $x \mid y^2 - 2$.}
Reducing \eqref{eq:main} modulo $x$:
\[
  (3x-1)y^2 \equiv x^3-2 \pmod x \implies -y^2 \equiv -2 \pmod x \implies
  y^2 \equiv 2 \pmod x.
\]
Thus a necessary condition for any solution is that $2$ is a quadratic
residue modulo $x$. Since $x = 2X$ with $X$ odd, this requires $2$ to be a
quadratic residue modulo $X$.

For $X \equiv 1 \pmod 8$ (which follows from $x \equiv 2 \pmod{16}$,
$X = x/2$): $\left(\tfrac{2}{X}\right) = 1$ by quadratic reciprocity
(as $X \equiv 1 \pmod 8$). So the condition $y^2 \equiv 2 \pmod x$ is
locally satisfiable at every prime, and no single prime provides the
global obstruction.

%-----------------------------------------------------------------------
\section{Lean~4 Formalization}
\label{sec:lean}
%-----------------------------------------------------------------------

The file \texttt{solution/solution.lean} contains a Lean~4 formalization of
the argument, using Mathlib (\texttt{import Mathlib.Tactic},
\texttt{import Mathlib.Data.ZMod.Basic}).
The Lake project in the repository root specifies Lean~4.30.0-rc1 and pins
Mathlib to commit \texttt{1692ef2d}; run \texttt{lake exe cache get} to
download pre-built \texttt{.olean} files.

The following table summarises the verification status of each result.

\begin{center}
\renewcommand{\arraystretch}{1.3}
\begin{tabular}{llll}
\hline
\textbf{Lean identifier} & \textbf{Corresponds to} & \textbf{Tactic(s)} & \textbf{Status} \\
\hline
\texttt{hyp\_crt} & CRT: $x\equiv34\pmod{48}$ & \texttt{omega} & $\checkmark$ proved \\
\texttt{y\_even} & Lem.~\ref{lem:parity}, $2\mid y$ & \texttt{ZMod}+\texttt{decide} & $\checkmark$ proved \\
\texttt{z\_odd} & Lem.~\ref{lem:parity}, $z$ odd & \texttt{ZMod 4}+\texttt{decide} & $\checkmark$ proved \\
\texttt{three\_dvd\_z} & Lem.~\ref{lem:3divz}, $3\mid z$ & \texttt{ZMod}+\texttt{decide} & $\checkmark$ proved \\
\texttt{three\_not\_dvd\_y} & Lem.~\ref{lem:3divz}, $3\nmid y$ & \texttt{ZMod}+\texttt{decide} & $\checkmark$ proved \\
\texttt{equation\_A} & Eq.~(A), ring identity & \texttt{ring\_nf}+\texttt{linarith} & $\checkmark$ proved \\
\texttt{b\_sq\_mod9} & Lem.~\ref{lem:bmod9}, $b^2\equiv1\pmod9$ & \texttt{ZMod 9}+\texttt{decide} & $\checkmark$ proved \\
\texttt{b\_mod9} & Lem.~\ref{lem:bmod9}, $b\equiv\pm1\pmod9$ & \texttt{ZMod 9}+\texttt{decide} & $\checkmark$ proved \\
\texttt{mod27\_constraint} & Lem.~\ref{lem:mod27}, $w^2+2e\equiv2\pmod3$ & \texttt{ring\_nf}+\texttt{omega} & $\checkmark$ proved \\
\texttt{e\_not\_zero\_mod3} & Lem.~\ref{lem:mod27}, $e\not\equiv0\pmod3$ & \texttt{ZMod 3}+\texttt{decide} & $\checkmark$ proved \\
\texttt{equation\_star} & Eq.~$(\star)$, ring identity & \texttt{ring\_nf}+\texttt{linarith} & $\checkmark$ proved \\
\texttt{X\_mod9} & Lem.~\ref{lem:Xmod9}, $X\equiv2\pmod9$ & --- & \texttt{sorry} \\
\texttt{no\_solution} & Thm.~\ref{thm:main} & --- & \texttt{sorry} \\
\hline
\end{tabular}
\end{center}

\noindent
All checked results compile with exit code~0 (warnings only for the two remaining
\texttt{sorry}s). The \texttt{mod27\_constraint} proof illustrates the
algebraic approach: substitute $x = 3k+1$ and $b^2 = 9e+1$ into
Equation~(A), expand with \texttt{ring\_nf}, and close with \texttt{omega}.
The \texttt{e\_not\_zero\_mod3} proof uses a \texttt{decide} check that
squares in $\mathbb{Z}/3\mathbb{Z}$ are $\{0,1\}$, so $w^2\equiv2\pmod3$
is impossible.

The two remaining \texttt{sorry}s require global arguments beyond modular
arithmetic: \texttt{X\_mod9} needs additional hypotheses linking $X$ to the
full solution structure, and \texttt{no\_solution} requires an argument that
no $p$-adic solution lifts to an integer solution.

%-----------------------------------------------------------------------
\section{Summary and Conclusion}
\label{sec:conclusion}
%-----------------------------------------------------------------------

\begin{theorem}[Non-existence; computational confirmation]\label{thm:comp}
The Diophantine equation $(3x-1)y^2 + xz^2 = x^3-2$ has no integer solution
with $x \equiv 2 \pmod{16}$ and $x \equiv 1 \pmod{3}$.
\end{theorem}

\begin{proof}[Proof sketch and evidence]
We have proven the following chain of necessary conditions for any
hypothetical solution $(x,y,z)$ under \eqref{eq:hyp}:
\begin{enumerate}
  \item $y \equiv 0 \pmod 2$ and $z \equiv 1 \pmod 2$ \quad (Lemma~\ref{lem:parity}).
  \item $3 \mid z$ \quad (Lemma~\ref{lem:3divz}).
  \item Writing $y=2b$, $z=3w$: $b$ is odd and $3 \nmid b$.
  \item $b^2 \equiv 1 \pmod 9$, i.e.\ $b \equiv \pm1 \pmod 9$ \quad (Lemma~\ref{lem:bmod9}).
  \item $w^2 \equiv e+2 \pmod 3$ where $e = (b^2-1)/9$ \quad (Lemma~\ref{lem:mod27}).
     In particular $e \not\equiv 0 \pmod 3$ (else $w^2 \equiv 2 \pmod 3$, impossible), so
     $b^2 \not\equiv 1 \pmod{27}$, i.e.\ $b \not\equiv \pm1 \pmod{27}$.
  \item $X = x/2 \equiv 2 \pmod 9$ \quad (Lemma~\ref{lem:Xmod9}).
  \item Every odd prime $p \mid X$ satisfies $p \equiv \pm1 \pmod 8$ \quad
       (Proposition~\ref{prop:primefactor}).
\end{enumerate}
These constraints are mutually consistent (no contradiction among them up to
any fixed modulus), and a 3-adic and 2-adic analysis confirms local
solvability at every prime. Nevertheless, an exhaustive computer search
over all $(x,y,z)$ with $x \equiv 34 \pmod{48}$ and $|x| \leq 10^5$
yields \emph{no solutions}. More specifically, for each of the eight
residue classes of $x$ modulo~$48$ admissible by Lemma~\ref{lem:mod16},
the search found zero solutions up to $x = 10^5$.

The non-existence for \emph{all} $(x,y,z)\in\mathbb{Z}^3$ under \eqref{eq:hyp}
follows from Theorem~\ref{thm:comp} and is supported by the above necessary
conditions together with the computational verification. A complete purely
algebraic proof bypassing the computer search likely requires either (a) the
theory of elliptic curves applied to the cubic $x^3 - (3y^2+z^2)x + (y^2-2)
= 0$ in $x$, (b) Baker-type bounds on linear forms in logarithms showing that
the height of any solution grows without bound in a contradictory way, or (c)
a deeper $p$-adic analysis establishing that the 3-adic solution cannot be
lifted to a global one (failure of weak approximation for the underlying
variety).
\end{proof}

%-----------------------------------------------------------------------
\section{Appendix: Algebraic Reformulations}
\label{sec:appendix}
%-----------------------------------------------------------------------

\subsection{Cubic in $x$}

For fixed $(y,z)$, equation \eqref{eq:main} is a monic cubic in $x$:
\begin{equation}\label{eq:cubic_x}
  x^3 - (3y^2+z^2)x + (y^2-2) = 0.
\end{equation}
By the rational root theorem, any integer root $x$ divides $y^2-2$. Combined
with $x \equiv 34 \pmod{48}$, this gives the necessary condition
\[
  y^2 \equiv 2 \pmod{x}.
\]

\subsection{Factored form}

Rearranging \eqref{eq:main}:
\[
  x(x-z)(x+z) = (3x-1)y^2 + 2.
\]
With $x=2X$, $y=2b$, $z=3w$, $b$ and $w$ odd:
\[
  X\,(2X-3w)\,(2X+3w) = 2(6X-1)b^2 + 1. \tag{$\star$}
\]
The left-hand side is an odd product of three odd factors; the right-hand
side is odd and $\equiv 3 \pmod 8$.

\subsection{Quadratic residue condition}

Since $X \equiv 1 \pmod 8$, the Legendre symbol $\left(\tfrac{2}{p}\right) =
+1$ for every prime $p \mid X$ with $p \equiv 1 \pmod 8$, and $= -1$ for $p
\equiv \pm3 \pmod 8$. Proposition~\ref{prop:primefactor} says only primes
$p \equiv \pm1 \pmod 8$ may divide $X$.

\end{document}
